% part: intuitionistic-logic
% chapter: propositions-as-types
% section: proofs-to-terms

\documentclass[../../../include/open-logic-section]{subfiles}

\begin{document}

\olfileid[pt]{int}{pty}{pt}

\olsection{Conversão de \usetoken{P}{derivation} em termos de prova}

O processo de converter !!{proof}s de dedução natural em termos de prova,
isto é, de traduzir !!{proof}s em~\Log{N2i} para !!{proof}s em~\Log{tN2i},
é extremamente direto. Basta escrever um termo de prova à esquerda de
cada !!{formula} que aparece à direita de um sequente na !!{derivation},
obtendo sequentes da forma $\Gamma \Sequent M : !A$. Diremos, então, que
$M$ \emph{testemunha}~$!A$ no contexto~$\Gamma$. O termo de prova a ser
escrito é inteiramente determinado pelas regras de \Log{tN2i}.

\begin{prop}
Toda !!{proof}~$\delta$ de $\Gamma \Sequent !A$ em \Log{N2i}
corresponde a !!a{proof}~$\delta'$ de $\Gamma \Sequent N : !A$ em
\Log{tN2i}. O termo de prova $N$ é unicamente determinado por~$\delta$.
\end{prop}

\begin{proof}
Procedemos por indução sobre a altura de~$\delta$. Se $\delta$ é o axioma
$x:!A \Sequent !A$, então $\delta'$ é o axioma~$x:!A \Sequent x:!A$. Se
$\delta$ termina em uma regra de inferência, a hipótese indutiva fornece
termos de prova para cada premissa e \Log{tN2i}-!!{proof}s de sequentes
que incluem esses termos de prova correspondentes a cada premissa.
Aplicamos a regra correspondente de \Log{tN2i}; o termo de prova
correspondente é determinado pela regra de inferência de \Log{tN2i}.
\end{proof}

\begin{ex}
  Considere a seguinte !!{proof} muito simples, em \Log{N1i}, de $(!A \land !B)
  \lif (!B \land !A)$:
  \[
  \AxiomC{$\Discharge{!A\land !B}{x}$}
  \RightLabel{\Elim\land[2]}
  \UnaryInfC{$!B$}
  \AxiomC{$\Discharge{!A\land !B}{x}$}
  \RightLabel{\Elim\land[1]}
  \UnaryInfC{$!A$}
  \RightLabel{\Intro\land}
  \BinaryInfC{$!B \land !A$}
  \DischargeRule{\Intro\lif}{x}
  \UnaryInfC{$(!A \land !B) \lif (!B \land !A)$}
  \DisplayProof
  \]
  Em \Log{N2i}, ela tem a seguinte forma:
  \[
  \Axiom$x : !A\land !B \fCenter !A\land !B$
  \RightLabel{\Elim\land[2]}
  \UnaryInf$x : !A\land !B \fCenter !B$
  \Axiom$x : !A\land !B \fCenter !A\land !B$
  \RightLabel{\Elim\land[1]}
  \UnaryInf$x : !A\land !B \fCenter !A$
  \RightLabel{\Intro\land}
  \BinaryInf$x : !A\land !B \fCenter !B \land !A$
  \RightLabel{\Intro\lif}
  \UnaryInf$\fCenter (!A \land !B) \lif (!B \land !A)$
  \DisplayProof
  \]
  Atribuímos termos de prova de cima para baixo. Os termos de prova dos
  axiomas são simplesmente a própria variável do contexto,~$x$. Os
  termos de prova associados a $\Elim\land[i]$ ficam, então, determinados
  como $\proj{2}{x}$ e~$\proj{1}{x}$. Ao aplicarmos \Intro{\land},
  combinamos esses dois termos de prova: $\pair{\proj{2}{x}}{\proj{1}{x}}$.
  Por fim, a regra \Intro{\lif} aplica uma abstração~$\lambd$, ligando~$x$
  e registrando que $x$ rotulava a hipótese~$!A \land !B$:
  $\lambd[\typeof{x}{!A \land
  !B}][\pair{\proj{2}{x}}{\proj{1}{x}}]$.
  A !!{proof} em \Log{tN2i} é, portanto:
  \[
  \Axiom$x : !A\land !B \fCenter x: !A\land !B$
  \RightLabel{\Elim\land[2]}
  \UnaryInf$x : !A\land !B \fCenter \proj{2}{x} : !B$
  \Axiom$x : !A\land !B \fCenter x: !A\land !B$
  \RightLabel{\Elim\land[1]}
  \UnaryInf$x : !A\land !B \fCenter \proj{1}{x} : !A$
  \RightLabel{\Intro\land}
  \BinaryInf$x : !A\land !B \fCenter \pair{\proj{2}{x}}{\proj{1}{x}} : 
    !B \land !A$
  \RightLabel{\Intro\lif}
  \UnaryInf$\fCenter \lambd[\typeof{x}{!A \land
  !B}][\pair{\proj{2}{x}}{\proj{1}{x}}] : (!A \land !B) \lif (!B \land !A)$
  \DisplayProof
  \]
\end{ex}


\end{document}
